\section{A differential-polynomial finite-part calculus}
\label{poly:calculus}

\subsection{Finite local data and the global normalization}

For $n\ge0$, put $b_n=H_n-\gamma$ and
$c_{r,n}=H_n^{(r)}+(-1)^r\zeta(r)$ for $r\ge2$. The familiar
Gamma recurrence and Taylor expansion give
\begin{equation}\label{poly:localpsi}
 v_n(\varepsilon):=\psi(-n+\varepsilon)
 =-\varepsilon^{-1}+b_n+
      \sum_{r\ge2}c_{r,n}\varepsilon^{r-1}.
\end{equation}
This series converges for $0<|\varepsilon|<1$; its differentiated
series give all polygammas. Let $M$ be an upper bound for the differential
weights of the monomials of $P$. Define
\begin{equation}\label{poly:poles}
 p_j(n)=[\varepsilon^{-j}]P(v_n,v_n',\ldots,v_n^{(R)}),\quad
 P_n(z)=\sum_{j=1}^M\frac{p_j(n)}{(z+n)^j},\quad
 \kappa_P=[\varepsilon^0]f_P(\varepsilon).
\end{equation}
All coefficients are finite polynomials in the stated harmonic numbers
and zeta values. To find the coefficients through the constant term,
one may truncate $v_n$ after $\varepsilon^{M-1}$ before differentiating.
For a homogeneous monomial of weight $M$, the regular coefficient of
harmonic order $r$ raises the exponent relative to its all-pole term
by $r$; hence orders above $M$ cannot contribute. Linearity handles
inhomogeneous polynomials.

\begin{theorem}[Normalized differential-polynomial expansion]
\label{poly:ML}
Away from the nonpositive integers,
\begin{equation}\label{poly:MLeq}
 f_P(z)=P_0(z)+\kappa_P+
              \sum_{n\ge1}\bigl(P_n(z)-P_n(0)\bigr).
\end{equation}
The series converges normally on compact subsets of its domain.
The entire remainder is exactly the displayed constant.
\end{theorem}

\begin{proof}
Each $p_j(n)$ grows at most as a fixed power of $1+\log(n+1)$:
only $H_n$ is unbounded, whereas the higher harmonic numbers are bounded.
On a fixed compact set,
$(n+z)^{-j}-n^{-j}=O(n^{-j-1})$. This proves normal convergence.
To determine the remainder, put $g(w)=f_P(w)-P_0(w)$ and integrate
$zg(w)/(w(w-z))$ on $|w|=N+1/2$.

The right-half-plane asymptotics give $\psi(w)=O(1+\log N)$ and
bounded fixed positive derivatives. On the left use
$\psi(w)=\psi(1-w)-\pi\cot(\pi w)$ and its derivatives.
The cotangent and each fixed derivative are uniformly bounded on these
circles. For $|\im w|\ge1$ this follows from their exponential
representations; for $|\im w|<1$ the circles remain a fixed positive
distance from the integers. Thus $g=O((1+\log N)^d)$ for some fixed
$d$, and the contour integral tends to zero. The residues at $z$ and
$0$ are $g(z)$ and $-\kappa_P$; a pole at $-n$ contributes
$-P_n(z)+P_n(0)$. The residue theorem gives the assertion.
\end{proof}

\subsection{The scalar projection controls both boundary obstructions}

Define polynomials in $L$ by
\begin{align}
 C_P(L)&=\frac1{2\pi}\int_{-\pi}^{\pi}\phi_P(L+\ii\theta)\dd\theta,
 & C_P&=C_P(0),\label{poly:C}\\
 R_P(L)&=-\frac{\phi_P(L+\ii\pi)-\phi_P(L-\ii\pi)}{2\ii\pi}.
 \label{poly:R}
\end{align}
The integrals are integrals of polynomials, so these are finite algebraic
formulas. If $\phi_P(Y)=\sum a_kY^k$, then
\begin{equation}\label{poly:Cexplicit}
 C_P=\sum_{m\ge0}a_{2m}\frac{(-1)^m\pi^{2m}}{2m+1}.
\end{equation}

\begin{theorem}[Scalar projection and residue cancellation]
\label{poly:filter}
For every fixed $P$, with $r_P(n)=p_1(n)$, there is an integer $d$ such that
\begin{align}
 \kappa_P-\sum_{n=1}^NP_n(0)&=C_P(\log N)+o(1),\label{poly:circle}\\
 r_P(n)&=R_P(\log n)+O((1+\log n)^d/n),\label{poly:rasym}\\
 r_P(n)-r_P(n-1)&=\frac{R_P'(\log n)}n
                      +O((1+\log n)^d/n^2).\label{poly:deltarasym}
\end{align}
In particular, a mixed polynomial has $C_P=R_P=0$.
\end{theorem}

\begin{proof}
The first identity follows by integrating $f_P(w)/w$ on the same
circles. Away from an arc of angular length $O(\log N/N)$ around the
negative real axis, reflection and the standard sectorial asymptotics
give $\psi(w)=\Log w+O(1/N)$ and
$\psi^{(r)}(w)=O(N^{-r})+O(N^{-2\pi})$ for fixed $r\ge1$.
On the exceptional arc all fixed polygammas are bounded and the digamma
is $O(\log N)$. Every mixed monomial consequently has angular average
$o(1)$. The scalar terms have the stated logarithmic average. Replacing
$\log(N+1/2)$ by $\log N$ changes the answer by $o(1)$.

For the residue, take the constant-tail specialization of the local
series: $b_n$ tends formally to $L$, and
$c_{r,n}$ tends to $(1+(-1)^r)\zeta(r)$. The resulting series is
$L+y$, where $y=-\pi\cot(\pi\varepsilon)$. Its derivatives satisfy
\[
 y'=y^2+\pi^2,
 \qquad y^{(r)}=Q_r(y),\quad
 Q_1(Y)=Y^2+\pi^2,\quad
 Q_{r+1}(Y)=(Y^2+\pi^2)Q_r'(Y).
\]
Every $Q_r$ with $r\ge1$ is divisible by $Y^2+\pi^2$.
Under the local change of variable $Y=y$, residue extraction is
\begin{equation}\label{poly:reschange}
 \Res_{\varepsilon=0}P(L+y,y',\ldots)
 =-[Y^{-1}]
 \frac{P(L+Y,Q_1(Y),Q_2(Y),\ldots)}{Y^2+\pi^2},
\end{equation}
where the coefficient is in the Laurent expansion at infinity.
The minus sign follows from $y\sim-1/\varepsilon$.
Every mixed term in the quotient is a polynomial and contributes zero.
For the scalar term, polynomial division leaves a remainder $AY+B$,
where
$A=(\phi_P(L+\ii\pi)-\phi_P(L-\ii\pi))/(2\ii\pi)$.
This proves the exact residue formula \eqref{poly:R}.

Finally, the harmonic asymptotic expansions through their $1/n$ terms
give $r_P(n)=R_P(\log n)+n^{-1}U(\log n)
+O((1+\log n)^d/n^2)$ for a polynomial $U$.
Differencing these explicit terms gives \eqref{poly:deltarasym}.
An unspecified $O(1/n)$ remainder is not being differentiated or
differenced to claim a stronger error estimate.
\end{proof}

\subsection{The complete spectral formula}

Write $\Delta p_j(n)=p_j(n)-p_j(n-1)$ and set
\begin{equation}\label{poly:JD}
 J_P=\sum_{n\ge1}\sum_{j=2}^M
          \frac{\Delta p_j(n)}{(j-1)n^{j-1}},\qquad
 D_P(s)=\sum_{n\ge1}\frac{\Delta p_1(n)}{n^s}\quad(\re s>0).
\end{equation}
The first sum is absolutely convergent. To see this directly, replace
$H_n^{(r)}$ by $H_{n-1}^{(r)}+n^{-r}$ in the finite pole polynomial;
each difference gains at least one factor $1/n$.

\begin{theorem}[Every differential-polynomial finite part]
\label{poly:bridge}
The Dirichlet series $D_P$ continues meromorphically to $\re s>-1$,
and
\begin{equation}\label{poly:bridgeeq}
 \boxed{\cI(P)=C_P+J_P+[s]D_P(s).}
\end{equation}
If $R_P'(L)=\sum a_jL^j$, an entirely ordinary-series form is
\begin{equation}\label{poly:ordinary}
 \cI(P)=C_P+J_P-\sum_j a_j\gamma_{j+1}
 -\sum_{n\ge1}
 \left\{\Delta p_1(n)-\frac{R_P'(\log n)}n\right\}\log n.
\end{equation}
The last series converges absolutely. If $\phi_P$ has degree $d\ge2$
and leading coefficient $a$, then $D_P$ has a pole of exact order
$d-1$ at zero, with leading term $-a\,d!s^{1-d}$. If
$\deg\phi_P\le1$, including $\phi_P=0$, it is holomorphic there.
\end{theorem}

\begin{proof}
The asymptotic in \eqref{poly:deltarasym} gives the normal decomposition
\[
 D_P(s)=\sum_j a_j(-1)^j\zeta^{(j)}(1+s)
 +\sum_{n\ge1}
 \left\{\Delta p_1(n)-R_P'(\log n)/n\right\}n^{-s}.
\]
Its second term is holomorphic on $\re s>-1$, including all derivatives
on compact subsets. By \eqref{poly:stieltjes}, its linear coefficient
is exactly the Stieltjes term and the last sum in
\eqref{poly:ordinary}. The leading term of $R_P'$ is
$-a d(d-1)L^{d-2}$, proving the pole classification.

It remains to establish the finite-part normalization. Integrate
\eqref{poly:MLeq}, first through a finite cutoff $N$. For $j\ge2$,
finite summation by parts gives
\[
 \sum_{n=1}^Np_j(n)(n^{1-j}-(n+1)^{1-j})
 =p_j(0)+\sum_{n=1}^N\frac{\Delta p_j(n)}{n^{j-1}}
   -\frac{p_j(N)}{(N+1)^{j-1}}.
\]
The first term cancels exactly the finite part of the pole at zero;
the last tends to zero. For simple poles the finite identity is
\[
 \sum_{n=1}^Np_1(n)\log(1+1/n)
 =p_1(N)\log(N+1)-\sum_{n=1}^N\Delta p_1(n)\log n.
\]
The boundary $p_1(N)\log(N+1)$ has logarithmic polynomial
$L R_P(L)$, with zero constant coefficient. The large-circle term
has constant $C_P$. The Stieltjes limit
$\sum_{n\le N}\log^j(n)/n
=\log^{j+1}N/(j+1)+\gamma_j+o(1)$ shows that the constant term of
the final logarithmic sum is $-[s]D_P(s)$. The subtracted series in
\eqref{poly:MLeq} is uniformly convergent after the origin pole is
removed, including on $[0,1]$, so this cutoff calculation computes the
finite part \eqref{poly:fp}. This proves both formulas.
\end{proof}

For a mixed $P$ the result simplifies to the ordinary convergent identity
\begin{equation}\label{poly:mixed}
 \boxed{\cI(P)=J_P+D_P'(0)
       =J_P-\sum_{n\ge1}\Delta p_1(n)\log n.}
\end{equation}
The disappearance of singular Laurent coordinates is thus a proved
consequence of the local cotangent algebra, not a fitted cancellation.
